\item \textbf{Problema de repaso.} Después de una colisión entre una gran nave espacial y un asteroide, un disco de cobre de 28.0 m de radio y 1.20 m de espesor, a una temperatura de $850^\circ\text{C}$, flota en el espacio, girando en torno a su eje de simetría con una rapidez angular de 250 rad/s. Mientras el disco radia luz infrarroja, su temperatura cae a $20.0^\circ\text{C}$. Sobre el disco no actúa momento de torsión externo. (a) Encuentre el cambio de energía cinética del disco. (b) Determine el cambio de energía interna del disco. (c) Obtenga la cantidad de energía que radia.
